% ---- pg-067 （书页 55） ----

\[
=E\begin{bmatrix}
(X_1-\mu_1)^2 & (X_1-\mu_1)(X_2-\mu_2)\\
(X_2-\mu_2)(X_1-\mu_1) & (X_2-\mu_2)^2
\end{bmatrix}
\]

\[
=\begin{bmatrix}
E(X_1-\mu_1)^2 & E(X_1-\mu_1)(X_2-\mu_2)\\
E(X_2-\mu_2)(X_1-\mu_1) & E(X_2-\mu_2)^2
\end{bmatrix}
\]

\[
=\begin{bmatrix}\sigma_{11} & \sigma_{12}\\\sigma_{21} & \sigma_{22}\end{bmatrix}=\begin{bmatrix}0.69 & -0.08\\-0.08 & 0.16\end{bmatrix}
\]

我们注意到离散型随机变量的均值、方差和协方差的计算涉及累加(如例 2.12 和例 2.13),
而对连续型随机变量类似的计算则涉及积分.

因为 $\sigma_{ik}=E(X_i-\mu_i)(X_k-\mu_k)=\sigma_{ki}$,可以方便地把\hyperref[eq:2-31]{式(2-31)}中的矩阵写成

\phantomsection\label{eq:2-32}
\begin{equation}
\bm{\Sigma}=E(\bm{X}-\bm{\mu})(\bm{X}-\bm{\mu})'=\begin{bmatrix}
\sigma_{11} & \sigma_{12} & \cdots & \sigma_{1p}\\
\sigma_{12} & \sigma_{22} & \cdots & \sigma_{2p}\\
\vdots & \vdots & & \vdots\\
\sigma_{1p} & \sigma_{2p} & \cdots & \sigma_{pp}
\end{bmatrix} \tag{2-32}
\end{equation}

$\bm{\mu}$ 和 $\bm{\Sigma}$ 可分别称为总体均值(向量)和总体方差-协方差矩阵(矩阵).

多元正态分布在均值向量 $\bm{\mu}$ 和方差-协方差矩阵 $\bm{\Sigma}$ 给定的情况下是完全确定的(见第 4 章),因而这些量在许多多元方法中起着重要作用.

把包含在方差 $\sigma_{ii}$ 中的信息,与包含在关联性度量,尤其正被称为总体相关系数 $\rho_{ik}$ 的关联性度量中的信息分离,常常能提供更多的信息. 相关系数 $\rho_{ik}$ 可用协方差 $\sigma_{ik}$ 和方差 $\sigma_{ii}$ 和 $\sigma_{kk}$ 定义为

\phantomsection\label{eq:2-33}
\begin{equation}
\rho_{ik}=\frac{\sigma_{ik}}{\sqrt{\sigma_{ii}}\,\sqrt{\sigma_{kk}}} \tag{2-33}
\end{equation}

相关系数度量随机变量 $X_i$ 和 $X_k$ 之间线性联系的程度(可参见文献[5]).

设总体相关矩阵是 $p\times p$ 对称阵

\phantomsection\label{eq:2-34}
\begin{equation}
\bm{\rho}=\begin{bmatrix}
\dfrac{\sigma_{11}}{\sqrt{\sigma_{11}}\,\sqrt{\sigma_{11}}} & \dfrac{\sigma_{12}}{\sqrt{\sigma_{11}}\,\sqrt{\sigma_{22}}} & \cdots & \dfrac{\sigma_{1p}}{\sqrt{\sigma_{11}}\,\sqrt{\sigma_{pp}}}\\[2ex]
\dfrac{\sigma_{12}}{\sqrt{\sigma_{11}}\,\sqrt{\sigma_{22}}} & \dfrac{\sigma_{22}}{\sqrt{\sigma_{22}}\,\sqrt{\sigma_{22}}} & \cdots & \dfrac{\sigma_{2p}}{\sqrt{\sigma_{22}}\,\sqrt{\sigma_{pp}}}\\[2ex]
\vdots & \vdots & & \vdots\\[1ex]
\dfrac{\sigma_{1p}}{\sqrt{\sigma_{11}}\,\sqrt{\sigma_{pp}}} & \dfrac{\sigma_{2p}}{\sqrt{\sigma_{22}}\,\sqrt{\sigma_{pp}}} & \cdots & \dfrac{\sigma_{pp}}{\sqrt{\sigma_{pp}}\,\sqrt{\sigma_{pp}}}
\end{bmatrix}
=\begin{bmatrix}
1 & \rho_{12} & \cdots & \rho_{1p}\\
\rho_{12} & 1 & \cdots & \rho_{2p}\\
\vdots & \vdots & & \vdots\\
\rho_{1p} & \rho_{2p} & \cdots & 1
\end{bmatrix} \tag{2-34}
\end{equation}

并设 $p\times p$ 标准差矩阵是

\phantomsection\label{eq:2-35}
\begin{equation}
\bm{V}^{1/2}=\begin{bmatrix}
\sqrt{\sigma_{11}} & 0 & \cdots & 0\\
0 & \sqrt{\sigma_{22}} & \cdots & 0\\
\vdots & \vdots & & \vdots\\
0 & 0 & \cdots & \sqrt{\sigma_{pp}}
\end{bmatrix} \tag{2-35}
\end{equation}
